% !PRINTED 186
\weberpage{186}
%<<CONT
\[
\left|
\begin{array}{l}
F''(x_1, x_1),\ F''(x_1, x_2)\ldots F''(x_1, x_m) \\
F''(x_2, x_1),\ F''(x_2, x_2)\ldots F''(x_2, x_m) \\
\cdots\cdots\cdots\cdots\cdots\cdots \\
F''(x_m, x_1),\ F'''(x_m, x_2)\ldots F''(x_m, x_m) % [sic] printed F''' in the second entry of the last row, inconsistent with F'' elsewhere
\end{array}
\right|
\]
and by $H'$ the corresponding determinant for the transformed
function $\Phi$, then by \S.\,56, (8)

\weq{6}{H' = r^2 H.}

$H$ is called the \gesperrt{Hessian determinant}\footnote{Named after \gesperrt{Hesse}, who made frequent use of this determinant in geometrical investigations, in particular in the theory of curves of the third order and in elimination problems, and who recognised its significance.} of the function
$F$. In both formulas (5) and (6), $r$ denotes the substitution
determinant.

\begin{center}
\S.\,60.\\
\gesperrt{Invariants and Covariants.}
\end{center}

In the linear transformation of homogeneous functions
there always occur certain functions of the coefficients which,
when one replaces the coefficients of the given form by the coefficients
of the transformed form, change in no other way than
that they take on a factor which is a power of the substitution
determinant. Such functions also occur in the simultaneous
transformation of a system of several forms, and they
depend on the coefficients of all the forms of the
system. Such functions are, for example, in a
system of $n$ linear forms the determinant of the system,
and in a quadratic form the determinant of this form.

These functions are called \gesperrt{invariants of the form}, or,
when they depend on several forms, \gesperrt{simultaneous invariants of the system}.

If, then, $I(a)$ is a function of the coefficients $a$ of a
form $F(x)$, and if $a'$ are the coefficients of the transformed
form $\Phi(x')$, then $I(a)$ is called an invariant if the identical relation
%CONT>>
